\begin{abstract}
    Este artículo se centra en cómo estructurar el análisis y diseño de una aplicación, para poder generar el código objetivo a través de modelos de razonamiento automático (LLMs). Pero esto solo cubre parte de la finalidad última de las herramientas generadas, que sería mas bien cómo estructurarlos para poder tener una trazabilidad completa desde el análisis a la implementación, y como automáticamente realizar las transformaciones de estos modelos en las 2 posibles direcciones; del análisis al codigo, y del código al análisis. Para ello, seguimos un enfoques de arquitectura dirigida por modelos (MDA), postulando que, centrándonos en un ámbito de programación muy concreto, en vez de ver a los requisitos como meros textos descriptivos de la funcionalidad, los debemos dotar de una estricta estructura en función del dominio específico en los que se mueven. Para ello, tomamos el ejemplo del análisis, diseño y generación de código para el desarrollo de sintetizadores de sonido software. Prescindiendo de modelar con UML, el cual, necesareamente estereotipado puede ser verboso en exceso, se propone que tanto el análisis como el diseño se estructuren, respetando sus distintos niveles de abstracción, en lenguajes de dominio específico (DSL) propios, concretos para el dominio genérico del tipo de aplicaciones a generar (en el ejemplo, síntesis de sonido), y con miras a que puedan ser facilmente entendidos tanto por una IA como por un analísta especializado en el dominio, y alejado de los detalles de implementación.
\end{abstract}
